\section{三阶段框架：Pre-training, Post-training, Inference-time}

\subsection{统一视角：能力形成的三层杠杆}
讲者把训练体系拆成三层杠杆：预训练决定底座容量与知识覆盖，后训练决定行为风格与任务偏好，推理时扩展决定单位请求可投入的搜索计算量。三者互补且存在明显的边际收益差异。

\begin{importantbox}{三阶段不是替代关系，而是分工关系}
\begin{itemize}
    \item \textbf{Pre-training}：提供广覆盖表征能力与知识压缩；
    \item \textbf{Post-training}：把能力定向到可用行为（对话、指令遵循、偏好一致）；
    \item \textbf{Inference-time}：在部署时动态投入额外计算，提高难任务成功率。
\end{itemize}
\end{importantbox}

\begin{figure}[H]
\centering
\includegraphics[width=0.9\textwidth]{frame-02-training-stack.jpg}
\caption{视频约 00:11:40：讲者展示从预训练到后训练再到测试时扩展的整体栈}
\end{figure}

\subsection{工程化分解：每层解决哪类瓶颈}
\begin{table}[H]
\centering
\renewcommand{\arraystretch}{1.2}
\begin{tabular}{p{2.6cm}p{3.2cm}p{4.0cm}p{3.2cm}}
\toprule
\textbf{阶段} & \textbf{主要目标} & \textbf{典型技术} & \textbf{常见瓶颈} \\
\midrule
Pre-training & 学表示、压缩知识、打基础能力 & 数据配比、优化器、学习率调度、架构细节 & 数据质量上限、算力成本高 \\
Post-training & 对齐行为、增强特定能力 & SFT、DPO/PPO、RLVR & 数据分布冲突、reward 偏差 \\
Inference-time & 提高难题成功率与鲁棒性 & Best-of-N、自一致性、检索迭代 & 延迟与成本上升、误判放大 \\
\bottomrule
\end{tabular}
\caption{三阶段训练框架及其瓶颈分工}
\end{table}

\begin{knowledgebox}{为什么课程里要反复强调 ``组合''}
单一阶段很难独立解决推理问题。比如只做后训练会受限于底座能力，只做推理时扩展会受限于候选质量。因此更现实的路径是按预算做 ``可叠加组合''，而不是寻找万能算法。
\end{knowledgebox}

\subsection{本章小结}
三阶段框架提供了分析推理模型的最小完备视角：先看底座，再看对齐，最后看部署时计算分配。


